\section*{Capítulo \ref{ch:g9:elements-universe-earth} --- Los elementos en el Universo y en la Tierra}

\begin{solution}{exo:g9:elements-universe-earth:1}
El hidrógeno (alrededor del \qty{74}{\percent}) y el helio (alrededor
del \qty{25}{\percent}).
\end{solution}

\begin{solution}{exo:g9:elements-universe-earth:2}
El oxígeno en los dos casos.
\end{solution}

\begin{solution}{exo:g9:elements-universe-earth:3}
En el interior de las estrellas, y en sus explosiones.
\end{solution}

\begin{solution}{exo:g9:elements-universe-earth:4}
Alrededor del \qty{3.5}{\percent}.
\end{solution}

\begin{solution}{exo:g9:elements-universe-earth:5}
$46.6 + 27.7 = \qty{74.3}{\percent}$: alrededor de tres cuartas partes.
\end{solution}

\begin{solution}{exo:g9:elements-universe-earth:6}
Son gases muy ligeros: escaparon de la Tierra joven hacia el espacio (y
el helio no forma compuestos que hubieran podido retenerlo en las
rocas).
\end{solution}

\begin{solution}{exo:g9:elements-universe-earth:7}
$0.23 \times 70 = \qty{16.1}{kg}$.
\end{solution}

\begin{solution}{exo:g9:elements-universe-earth:8}
Los \omterm{def:g7:atoms-and-molecules:atom}{átomos} de oxígeno son 16 veces más pesados que los de hidrógeno:
aunque haya menos \omterm{def:g7:atoms-and-molecules:atom}{átomos} de oxígeno, pesan más.
\end{solution}

\begin{solution}{exo:g9:elements-universe-earth:9}
Alrededor de $0.081 \times 1000 = \qty{81}{kg}$ de aluminio.
\end{solution}

\begin{solution}{exo:g9:elements-universe-earth:10}
El oxígeno y el calcio (y, fuera de los ocho elementos principales de
la corteza, el hidrógeno, el carbono y el fósforo están presentes en la
corteza en pequeñas cantidades).
\end{solution}

\begin{solution}{exo:g9:elements-universe-earth:11}
$40 \times 1.67 \times 10^{-27} = \qty{6.68e-26}{kg}$;
$1.06 / 6.68 \times 10^{-26} \approx 1.6 \times 10^{25}$ \omterm{def:g7:atoms-and-molecules:atom}{átomos} de calcio.
\end{solution}

\begin{solution}{exo:g9:elements-universe-earth:12}
Los \omterm{def:g5:irreversible-changes:chemical-change}{cambios químicos} no crean ni destruyen \omterm{def:g7:atoms-and-molecules:atom}{átomos}: solo los hacen pasar
de un compuesto a otro (roca, gas, azúcar, ser vivo). Los \omterm{def:g7:atoms-and-molecules:atom}{átomos} de
carbono de la Tierra se usan una y otra vez.
\end{solution}

\begin{solution}{pb:g9:elements-universe-earth:1}
\textbf{1.} $0.61 \times 70 = \qty{42.7}{kg}$.

\textbf{2.} Carbono: $0.23 \times 70 = \qty{16.1}{kg}$; hidrógeno:
$0.10 \times 70 = \qty{7.0}{kg}$.

\textbf{3.} $61 + 23 + 10 = \qty{94}{\percent}$.

\textbf{4.} El agua, \ce{H2O}.

\textbf{5.} $16 \times 1.67 \times 10^{-27} = \qty{2.67e-26}{kg}$.

\textbf{6.} Carbono: $12 \times 1.67 \times 10^{-27} = \qty{2.00e-26}{kg}$;
hidrógeno: \qty{1.67e-27}{kg}.

\textbf{7.} 16.

\textbf{8.} Un \omterm{def:g9:inside-the-atom:nucleus}{electrón} es unas 1836 veces más ligero que un \omterm{def:g9:inside-the-atom:proton}{nucleón}:
los \omterm{def:g9:inside-the-atom:nucleus}{electrones} suponen menos de una milésima parte de la masa.

\textbf{9.} $7.0 / 1.67 \times 10^{-27} \approx 4.2 \times 10^{27}$
\omterm{def:g7:atoms-and-molecules:atom}{átomos} de hidrógeno.

\textbf{10.} $16.1 / 2.00 \times 10^{-26} \approx 8.0 \times 10^{26}$
\omterm{def:g7:atoms-and-molecules:atom}{átomos} de carbono.

\textbf{11.} $42.7 / 2.67 \times 10^{-26} \approx 1.60 \times 10^{27}$
\omterm{def:g7:atoms-and-molecules:atom}{átomos} de oxígeno.

\textbf{12.} Nuestro recuento, $1.60 \times 10^{27}$, concuerda con la
estimación $1.61 \times 10^{27}$: un cuerpo de \qty{70}{kg} contiene
unos $1.6 \times
10^{27}$ \omterm{def:g7:atoms-and-molecules:atom}{átomos} de oxígeno.
\end{solution}
